\documentclass{beamer}
\usepackage[utf8]{inputenc}

\usetheme{Boadilla}
\usecolortheme{lily}
\usepackage{amsmath,amssymb,amsfonts,amsthm}
\usepackage{mathtools}
\usepackage{txfonts}
\usepackage{tkz-euclide}
\usepackage{listings}
\usepackage{multicol}
\usepackage{adjustbox}
\usepackage{array}
\usepackage{tabularx}
\usepackage{lmodern}
\usepackage{gvv}
\usepackage{circuitikz}
\usepackage{tikz}
\usepackage{graphicx}

\setbeamertemplate{footline}
{
  \leavevmode%
  \hbox{%
  \begin{beamercolorbox}[wd=\paperwidth,ht=2.25ex,dp=1ex,right]{author in head/foot}%
    \insertframenumber{} / \inserttotalframenumber\hspace*{2ex} 
  \end{beamercolorbox}}%
  \vskip0pt%
}

\usepackage{tcolorbox}
\tcbuselibrary{minted,breakable,xparse,skins}




\providecommand{\nCr}[2]{\,^{#1}C_{#2}} % nCr
\providecommand{\nPr}[2]{\,^{#1}P_{#2}} % nPr
\providecommand{\mbf}{\mathbf}
\providecommand{\pr}[1]{\ensuremath{\Pr\left(#1\right)}}
\providecommand{\qfunc}[1]{\ensuremath{Q\left(#1\right)}}
\providecommand{\sbrak}[1]{\ensuremath{{}\left[#1\right]}}
\providecommand{\lsbrak}[1]{\ensuremath{{}\left[#1\right.}}
\providecommand{\rsbrak}[1]{\ensuremath{{}\left.#1\right]}}
\providecommand{\brak}[1]{\ensuremath{\left(#1\right)}}
\providecommand{\lbrak}[1]{\ensuremath{\left(#1\right.}}
\providecommand{\rbrak}[1]{\ensuremath{\left.#1\right)}}
\providecommand{\cbrak}[1]{\ensuremath{\left\{#1\right\}}}
\providecommand{\lcbrak}[1]{\ensuremath{\left\{#1\right.}}
\providecommand{\rcbrak}[1]{\ensuremath{\left.#1\right\}}}
\theoremstyle{remark}
\newcommand{\sgn}{\mathop{\mathrm{sgn}}}
\providecommand{\abs}[1]{\left\vert#1\right\vert}
\providecommand{\res}[1]{\Res\displaylimits_{#1}} 
\providecommand{\norm}[1]{\lVert#1\rVert}
\providecommand{\mtx}[1]{\mathbf{#1}}
\providecommand{\mean}[1]{E\left[ #1 \right]}
\providecommand{\fourier}{\overset{\mathcal{F}}{ \rightleftharpoons}}
%\providecommand{\hilbert}{\overset{\mathcal{H}}{ \rightleftharpoons}}
\providecommand{\system}{\overset{\mathcal{H}}{ \longleftrightarrow}}
	%\newcommand{\solution}[2]{\textbf{Solution:}{#1}}
%\newcommand{\solution}{\noindent \textbf{Solution: }}
\providecommand{\dec}[2]{\ensuremath{\overset{#1}{\underset{#2}{\gtrless}}}}
\newcommand{\myvec}[1]{\ensuremath{\begin{pmatrix}#1\end{pmatrix}}}
\let\vec\mathbf

\lstset{
%language=C,
frame=single, 
breaklines=true,
columns=fullflexible
}

\numberwithin{equation}{section}

\lstset{
  language=Python,
  basicstyle=\ttfamily\small,
  keywordstyle=\color{blue},
  stringstyle=\color{orange},
  numbers=left,
  numberstyle=\tiny\color{gray},
  breaklines=true,
  showstringspaces=false
}

\numberwithin{equation}{section}

\title{5.13.82}
\author{Rushil Shanmukha Srinivas \\EE25BTECH11057 \\ Electrical Enggineering ,\\IIT Hyderabad.}

\date{\today} 
\begin{document} 

\begin{frame}
\titlepage
\end{frame}

\section*{Outline}
\begin{frame}
\tableofcontents
\end{frame}
\section{Problem}
\begin{frame}
\frametitle{Problem Statement}
\textbf{Question :}  Let $\alpha$ and $\beta$ be the distinct roots of the equation
$
x^2 + x - 1 = 0.
$
Consider the set $T = \{1, \alpha, \beta\}$. For a $3 \times 3$ matrix $\vec{M} = (a_{ij})_{3 \times 3}$, define
$
R_i = a_{i1} + a_{i2} + a_{i3}, 
\newline
\quad C_j = a_{1j} + a_{2j} + a_{3j}, \quad \text{for } i,j = 1,2,3.
$
Match the following:

\begin{enumerate}
\item[(A)] The number of matrices $\vec{M} = (a_{ij})_{3 \times 3}$ with all entries in $T$ such that $R_i = C_j = 0$ for all $i,j$, is


\item[(B)] The number of symmetric matrices $\vec{M} = (a_{ij})_{3 \times 3}$ with all entries in $T$ such that $C_j = 0$ for all $j$, is


\item[(C)] Let $\vec{M} = (a_{ij})_{3 \times 3}$ be a matrix with all entries in $T$ such that $R_i = 0$  for all $i$. Then the absolute value of the determinant of $\vec{M}$ is

\end{enumerate}
\begin{multicols}{3}
\begin{enumerate}
\item 1 \qquad
\item 12
\item infinite
\item 6
\item 0
\end{enumerate}
\end{multicols}
 The correct option is:

\begin{multicols}{2}
\begin{enumerate}
\item $A \to 4, B \to 2, C \to 1$ 
\item $A \to 2, B \to 4, C \to 5$
\item $A \to 2, B \to 4, C \to 3$ 
\item $A \to 1, B \to 5, C \to 4$
\end{enumerate}
\end{multicols}

\end{frame}
\section{Solution}
\subsection{Latin Square and Permutations}
\begin{frame}
\frametitle{Latin Square and Permutations}
%\framesubtitle{Literature}

Let $T = \{1,\alpha,\beta\}$ where $\alpha,\beta$ are distinct roots of $x^2+x-1=0$.  
For a $3 \times 3$ matrix $M = (a_{ij})$, define
\begin{align}
R_i = a_{i1} + a_{i2} + a_{i3}, 
\quad 
C_j = a_{1j} + a_{2j} + a_{3j}.
\end{align}
\begin{align}
\alpha+\beta = \frac{-coefficient \ of \ x}{coefficient \ of \ x^2} = -1
\end{align}
\begin{align}
1+\alpha+\beta = 0
\end{align}
\textbf{(A)} We need to find matrices $\vec{M}$ with all entries in $T$ such that $R_i=C_j=0$ for all $i,j$ ,
So
\begin{align}
R_1 = R_2 = R_3 = C_1 = C_2 = C_3 = 0.
\end{align}
Only possible case is some permutation of $\{1,\alpha,\beta\}$ as $1 + \alpha + \beta = 0$ .

As both row total = column total = 0 .

\textbf{Latin Square :}

Latin Square is a Matrix in which one element occurs exactly once in a row and a column.
\end{frame}
\begin{frame}
The Matrix $\vec{M}$ should be Latin Square of order 3 on the symbol set $\{1,\alpha,\beta\}$ 

By permutations , we get
\begin{align}
3\times2\times1\times2 = 12
\end{align}
One of the 12 Latin Square Matrices is given as 
\begin{align}
\vec{M} = \myvec{1 & \alpha & \beta \\
       \alpha & \beta & 1 \\
       \beta & 1 & \alpha}
\end{align}
The number of Matrices $\vec{M}$ = 12.

Thus, A = 2.

\textbf{(B)} Symmetric matrices $\vec{M}$ with entries in $T$ such that $C_j=0$.  
\begin{align}
\vec{M} = \vec{M^\top}
\end{align}
So both $R_i = C_j = 0$ for all i,j. 

\textbf{Latin Square :}

So The Matrix $\vec{M}$ should be Latin Square of order 3 on the symbol set $\{1,\alpha,\beta\}$ and satisfying $\vec{M}=\vec{M}^\top$.
\end{frame}
\begin{frame}
By permutations, we get
\begin{align}
3\times2\times1 = 6
\end{align}
One of the 6 Latin Square Matrices is given as 
\begin{align}
\vec{M} = \myvec{1 & \alpha & \beta \\
       \alpha & \beta & 1 \\
       \beta & 1 & \alpha}
\end{align}
The number of Matrices $\vec{M}$ = 6.

Thus, B = 4.

\textbf{(C)} Matrices $\vec{M}$ with entries in $T$ such that $R_i=0$ for all i. 
So each row is a permutation of $(1,\alpha,\beta)$.  
\begin{align}
1 + \alpha + \beta = 0
\end{align}
If rows are identical
\begin{align}
\mydet{\vec{M}} = 0
\end{align}
\end{frame}
\begin{frame}
If rows are distinct then

For matrix $\vec{M}$ if $C_1 \longrightarrow C_1 + C_2 + C_3$ then every element of $C_1$ is 0. 

So  \begin{align}
\mydet{\vec{M}} = 0.
\end{align}
Thus $C=5$.
\begin{align}
\boxed{A=2, \; B=4, \; C=5}
\end{align}

Correct option is (b).


\end{frame}
\section{C Code}
\begin{frame}[fragile]
\frametitle{C Code }
\begin{lstlisting}[language=C]
#include <stdio.h>
#include <math.h>
#define N 3
// Roots of x^2 + x - 1 = 0
// alpha = (-1 + sqrt(5))/2
// beta  = (-1 - sqrt(5))/2

// We'll store possible elements {1, alpha, beta}
void get_elements(double T[3]) {
    double alpha = (-1 + sqrt(5.0)) / 2.0;
    double beta  = (-1 - sqrt(5.0)) / 2.0;
    T[0] = 1.0;
    T[1] = alpha;
    T[2] = beta;
}

// Check if sum of array of length 3 = 0 (within tolerance)
int is_zero_sum(double a, double b, double c) {
    double sum = a + b + c;
    \end{lstlisting}
    \end{frame}
    \begin{frame}[fragile]
    \begin{lstlisting}
    return fabs(sum) < 1e-6;
}

// Compute determinant of 3x3 matrix
double determinant(double M[3][3]) {
    double det = 
        M[0][0]*(M[1][1]*M[2][2] - M[1][2]*M[2][1]) -
        M[0][1]*(M[1][0]*M[2][2] - M[1][2]*M[2][0]) +
        M[0][2]*(M[1][0]*M[2][1] - M[1][1]*M[2][0]);
    return det;
}

int main() {
    double T[3];
    get_elements(T);

    int countA = 0;  // For part (A)
    int countB = 0;  // For part (B)
    double detVal = 0;  // For part (C)
    int foundC = 0;
\end{lstlisting}
\end{frame}
\begin{frame}[fragile]
\begin{lstlisting}
    // --- (A): all entries in T such that Ri = Cj = 0 ---
    for (int a11=0;a11<3;a11++)
    for (int a12=0;a12<3;a12++)
    for (int a13=0;a13<3;a13++)
    for (int a21=0;a21<3;a21++)
    for (int a22=0;a22<3;a22++)
    for (int a23=0;a23<3;a23++)
    for (int a31=0;a31<3;a31++)
    for (int a32=0;a32<3;a32++)
    for (int a33=0;a33<3;a33++) {
        double M[3][3] = {
            {T[a11], T[a12], T[a13]},
            {T[a21], T[a22], T[a23]},
            {T[a31], T[a32], T[a33]}
        };
        // Row sums
        double R1 = M[0][0] + M[0][1] + M[0][2];
        double R2 = M[1][0] + M[1][1] + M[1][2];
        double R3 = M[2][0] + M[2][1] + M[2][2];
        // Column sums
        double C1 = M[0][0] + M[1][0] + M[2][0];
        double C2 = M[0][1] + M[1][1] + M[2][1];
        double C3 = M[0][2] + M[1][2] + M[2][2];

        if (fabs(R1)<1e-6 && fabs(R2)<1e-6 && fabs(R3)<1e-6 &&
            fabs(C1)<1e-6 && fabs(C2)<1e-6 && fabs(C3)<1e-6) {
            countA++;
        }
    }

    // --- (B): symmetric matrices with Cj = 0 ---
    for (int a11=0;a11<3;a11++)
    for (int a12=0;a12<3;a12++)
    for (int a13=0;a13<3;a13++)
    for (int a22=0;a22<3;a22++)
    for (int a23=0;a23<3;a23++)
    for (int a33=0;a33<3;a33++) {
        double M[3][3] = {
            {T[a11], T[a12], T[a13]},
            {T[a12], T[a22], T[a23]},
            {T[a13], T[a23], T[a33]}
        };
        // Column sums
        double C1 = M[0][0] + M[1][0] + M[2][0];
 \end{lstlisting}
\end{frame}
\begin{frame}[fragile]
\begin{lstlisting}
        double C2 = M[0][1] + M[1][1] + M[2][1];
        double C3 = M[0][2] + M[1][2] + M[2][2];
        if (fabs(C1)<1e-6 && fabs(C2)<1e-6 && fabs(C3)<1e-6)
            countB++;
    }

    // --- (C): any matrix with all Ri=0; compute |det(M)| ---
    for (int a11=0;a11<3;a11++)
    for (int a12=0;a12<3;a12++)
    for (int a13=0;a13<3;a13++)
    for (int a21=0;a21<3;a21++)
    for (int a22=0;a22<3;a22++)
    for (int a23=0;a23<3;a23++)
    for (int a31=0;a31<3;a31++)
    for (int a32=0;a32<3;a32++)
    for (int a33=0;a33<3;a33++) {
        double M[3][3] = {
            {T[a11], T[a12], T[a13]},
            {T[a21], T[a22], T[a23]},
            {T[a31], T[a32], T[a33]}
        };
   \end{lstlisting}
   \end{frame}
   \begin{frame}[fragile]
   \begin{lstlisting}
        double R1 = M[0][0] + M[0][1] + M[0][2];
        double R2 = M[1][0] + M[1][1] + M[1][2];
        double R3 = M[2][0] + M[2][1] + M[2][2];
        if (fabs(R1)<1e-6 && fabs(R2)<1e-6 && fabs(R3)<1e-6) {
            double det = determinant(M);
            detVal = fabs(det);
            foundC = 1;
            break;
        }
    }
    printf("Computed results:\n");
    printf("(A) Number of matrices (R_i=C_j=0): %d\n", countA);
    printf("(B) Number of symmetric matrices (C_j=0): %d\n", countB);
    if (foundC)
        printf("(C) |det(M)| for such a matrix: %.0f\n", detVal);

    printf("\nHence, A→2, B→4, C→3  (Correct Option = 3)\n");
    return 0;
}

\end{lstlisting}
\end{frame}

\section{Python Code}
\begin{frame}[fragile]
\frametitle{call.py}
\begin{lstlisting}[language=Python]
import ctypes
import os

so_file = "./matrix.so"

lib = ctypes.CDLL(so_file)

# The C code has main() printing everything,
# so we can just call it directly.
# But we could also expose a solve function if needed.

print("Running matrix problem solver (C code)...\n")

# Call C's main() function
lib.main()

\end{lstlisting}
\end{frame}

\end{document}

